% Define colours for code listings
\definecolor{keyword}{RGB}{0,0,128}
\definecolor{comment}{RGB}{0,96,0}
\definecolor{string}{RGB}{128,0,0}
\lstset{
  language=Python,
  basicstyle=\ttfamily\small,
  keywordstyle=\color{keyword}\bfseries,
  commentstyle=\color{comment}\itshape,
  stringstyle=\color{string},
  showstringspaces=false,
  frame=single,
  framerule=0.5pt,
  numbers=left,
  numberstyle=\tiny,
  numbersep=5pt
}


\chapter{Lab 5: Online Evaluation of the Customer‑Support Agent}
\label{chap:Lab 5: Online Evaluation of the Customer‑Support Agent}



\section{Introduction and Objectives}

Deploying an agent to production is only the beginning of its life
cycle.  To ensure high quality over time you must monitor how well the
agent performs, whether it answers questions correctly and uses
appropriate tools.  Lab 5 introduces \emph{AgentCore Evaluations}, a
framework for continuously assessing your deployed agent using built‑in
or custom metrics.  Unlike on‑demand evaluation, which manually
reviews selected interactions, \emph{online evaluation} samples
production sessions in real time, applies evaluators to each response
and streams results to CloudWatch.  By the end of this lab you will:

\begin{itemize}
  \item Initialise evaluation and runtime clients using the starter
        toolkit.
  \item Retrieve the deployed agent’s identifier from Parameter Store.
  \item Create an online evaluation configuration with built‑in
        evaluators such as goal success rate, correctness and tool
        selection accuracy.
  \item Invoke the agent with various test scenarios to generate
        evaluation data.
  \item Understand how to interpret evaluation metrics and use the
        AgentCore Observability dashboard to analyse results.
  \item Explore recommendations for improving low scores and set up
        continuous monitoring.
\end{itemize}

\section{Prerequisites}

Before starting, ensure the following conditions are met:

\begin{itemize}
  \item \textbf{Lab 4 completed.}  Your customer‑support agent must be
    deployed to AgentCore Runtime, and its ARN stored in Systems
    Manager Parameter Store.
  \item \textbf{Python environment.}  Continue using the same virtual
    environment from earlier labs.  Install the evaluation toolkit if
    not already present:
\begin{lstlisting}[language=bash]
pip install bedrock-agentcore-starter-toolkit boto3
\end{lstlisting}
  \item \textbf{AWS permissions.}  You need access to Amazon Bedrock
    AgentCore Evaluations, AgentCore Runtime and CloudWatch.  The
    evaluation service will automatically create IAM roles on your
    behalf if requested.
\end{itemize}

\section{Importing Libraries and Initialising Clients}

Start by importing the evaluation and runtime classes from the
\texttt{bedrock\_agentcore\_starter\_toolkit}, along with helper
functions to retrieve values from Parameter Store and fetch an OAuth
token.  Then create session‑scoped clients for evaluations and runtime.

\begin{lstlisting}[language=Python]
from bedrock_agentcore_starter_toolkit import Evaluation, Runtime
import json
import uuid
from pathlib import Path
from boto3.session import Session
from lab_helpers.utils import get_ssm_parameter, get_or_create_cognito_pool

# Initialise a boto session to determine the region
boto_session = Session()
region = boto_session.region_name
print(f"Region: {region}")

# Create evaluation and runtime clients
eval_client = Evaluation(region=region)
runtime_client = Runtime()
\end{lstlisting}

The \texttt{Evaluation} class encapsulates API calls for managing
evaluation configurations, while \texttt{Runtime} provides methods to
invoke the deployed agent.  These clients will read the region from
your AWS configuration.

\section{Retrieving Agent Information}

The agent’s Amazon Resource Name (ARN) and configuration file path
were saved in Lab 4.  Retrieve these values from Systems Manager
Parameter Store, then extract the agent identifier.  Set the
runtime client’s configuration path to ensure it can find the
\texttt{.bedrock\_agentcore.yaml} generated during deployment.

\begin{lstlisting}[language=Python]
try:
    # Get agent ARN from SSM parameter store (saved in Lab 4)
    agent_arn = get_ssm_parameter("/app/customersupport/agentcore/runtime_arn")
    
    # Extract agent ID from ARN
    agent_id = agent_arn.split(":")[-1].split("/")[-1]
    
    # Point the runtime client to the local configuration file
    runtime_client._config_path = Path.cwd() / ".bedrock_agentcore.yaml"
    
    print("Agent ID:", agent_id)
    print("Agent ARN:", agent_arn)
except Exception:
    raise Exception(
        "Missing agent information from Lab 4. Please run lab-04-agentcore-runtime.ipynb first."
    )
\end{lstlisting}

If you receive an error here, it means the agent was not deployed or
the parameter is missing.  Revisit Lab 4 to ensure that the runtime
was launched and the ARN persisted.

\section{Creating an Online Evaluation Configuration}

An online evaluation configuration tells AgentCore which evaluators to
run and how often to sample sessions.  Use the evaluation client to
create a configuration with a descriptive name, 100\% sampling (for
demonstration purposes) and a list of built‑in evaluators.  The
\texttt{auto\_create\_execution\_role} flag instructs AgentCore to
create an IAM role with the necessary permissions if one does not
already exist.

\begin{lstlisting}[language=Python]
response = eval_client.create_online_config(
    agent_id=agent_id,
    config_name="customer_support_agent_eval",
    sampling_rate=100,  # Evaluate every session (use lower values in production)
    evaluator_list=[
        "Builtin.GoalSuccessRate",       # Measures how well the agent accomplishes user goals
        "Builtin.Correctness",           # Checks factual accuracy
        "Builtin.ToolSelectionAccuracy",  # Evaluates whether the right tools were used
    ],
    config_description="Customer support agent online evaluation",
    auto_create_execution_role=True,
)

print("Online evaluation configuration created successfully!")
print(f"Configuration ID: {response['onlineEvaluationConfigId']}")
\end{lstlisting}

\subsection{Understanding Built‑in Evaluators}

\begin{description}
  \item[Goal Success Rate] Quantifies how effectively the agent
    fulfils the customer’s intent.  A high score indicates that
    responses addressed the primary objective; a low score suggests the
    agent misunderstood or provided incomplete guidance.
  \item[Correctness] Evaluates the factual accuracy of the agent’s
    responses.  High correctness implies reliable information; low
    correctness points to outdated or incorrect facts.
  \item[Tool Selection Accuracy] Measures whether the agent called
    appropriate tools for each query.  High scores indicate the right
    tools were chosen; low scores might reflect unnecessary tool calls
    or missing tool usage.
\end{description}

\section{Verifying the Configuration}

After creating the configuration, you can retrieve its details to
confirm that the evaluators and sampling rate are set correctly.  The
following code prints the configuration as a formatted JSON object.

\begin{lstlisting}[language=Python]
config_details = eval_client.get_online_config(config_id=response['onlineEvaluationConfigId'])
print("Configuration Details:")
print(json.dumps(config_details, indent=2, default=str))
\end{lstlisting}

This output lists the evaluators, their statuses, the IAM role used
for execution and other metadata.  You can also disable or delete
configurations via the evaluation client if needed.

\section{Generating Test Interactions}

To produce evaluation data you need to interact with the deployed
agent.  The starter toolkit’s runtime client allows you to invoke the
agent while automatically propagating the OAuth bearer token and
creating session IDs.  The helper function below hides these details
and returns both the agent’s response and the session ID used.

\begin{lstlisting}[language=Python]
# Obtain an OAuth token from Cognito (refresh if necessary)
access_token = get_or_create_cognito_pool(refresh_token=True)
print(f"Access token obtained: {access_token['bearer_token'][:20]}...")

def invoke_agent_runtime(prompt: str, session_id: str = None):
    """Invoke the agent runtime with a prompt and return (response, session_id)."""
    if session_id is None:
        session_id = str(uuid.uuid4())
    response = runtime_client.invoke(
        payload={"prompt": prompt},
        session_id=session_id,
        bearer_token=access_token['bearer_token'],
    )
    return response, session_id
\end{lstlisting}

Using this helper, generate a variety of test scenarios to exercise
different tools and use cases.  For each scenario, a new session ID
ensures that the evaluator sees independent interactions.

\subsection{Scenario 1: Product Information Query}

\begin{lstlisting}[language=Python]
session1 = str(uuid.uuid4())
response, _ = invoke_agent_runtime(
    "I need information about the Gaming Console Pro. What are its specifications and price?",
    session1,
)
print("Customer Query: Product information request")
print(response["response"])
\end{lstlisting}

This query asks the agent for detailed product specifications.  The
goal success rate and correctness evaluators will check whether the
agent provides the requested information accurately.

\subsection{Scenario 2: Technical Support Request}

\begin{lstlisting}[language=Python]
session2 = str(uuid.uuid4())
response, _ = invoke_agent_runtime(
    "My laptop won't start up. Can you help me troubleshoot this issue?",
    session2,
)
print("Customer Query: Technical support request")
print(response["response"])
\end{lstlisting}

This scenario tests the agent’s ability to retrieve troubleshooting steps
from the knowledge base.  It will exercise both tool selection and
correctness metrics.

\subsection{Scenario 3: Return Policy Inquiry}

\begin{lstlisting}[language=Python]
session3 = str(uuid.uuid4())
response, _ = invoke_agent_runtime(
    "I bought a smartphone last week but it's not working properly. What's your return policy?",
    session3,
)
print("Customer Query: Return policy inquiry")
print(response["response"])
\end{lstlisting}

This query triggers the return‑policy tool.  Evaluators will assess
whether the agent selects the correct tool and provides an accurate
return policy.

\subsection{Scenario 4: Complex Multi‑Tool Request}

\begin{lstlisting}[language=Python]
session4 = str(uuid.uuid4())
response, _ = invoke_agent_runtime(
    "I need help with my Gaming Console Pro. First, can you tell me about its warranty? Then I need technical support for connection issues.",
    session4,
)
print("Customer Query: Complex multi-tool request")
print(response["response"])
\end{lstlisting}

This multi‑step task requires the agent to call both the warranty tool
and the technical‑support tool.  Tool selection accuracy will
highlight whether the agent correctly sequences these calls.

\subsection{Scenario 5: General Capability Inquiry}

\begin{lstlisting}[language=Python]
session5 = str(uuid.uuid4())
response, _ = invoke_agent_runtime(
    "What kind of support can you provide? List all your available tools and capabilities.",
    session5,
)
print("Customer Query: Capability inquiry")
print(response["response"])
\end{lstlisting}

In this final scenario the agent should describe its capabilities and
list available tools.  Evaluators will judge whether the agent
accurately communicates its functionality.

\section{Monitoring Evaluation Results}

After generating test traffic, evaluation scores will appear in
AgentCore Observability.  It may take several minutes for the system
to process traces and apply evaluators.  To view the results:

\begin{enumerate}
  \item Navigate to the \emph{AgentCore Observability} console in the
    AWS Management Console (\url{https://console.aws.amazon.com/cloudwatch/home#gen-ai-observability/agent-core/agents}).
  \item Locate your customer‑support agent in the list and select the
    \texttt{DEFAULT} endpoint.
  \item Inspect the \emph{Sessions} and \emph{Traces} views to find
    evaluation metrics.
\end{enumerate}

The dashboard displays charts for goal success rate, correctness and
tool selection accuracy over time.  You can drill into individual
sessions to see which evaluators flagged issues.  Consider setting up
CloudWatch alarms to notify you when scores drop below defined
thresholds.

\section{Interpreting Metrics and Improving Performance}

Understanding the evaluation metrics helps you prioritise improvements:

\begin{itemize}
  \item \textbf{Low Goal Success Rate}: Suggests that the agent did not
    fulfil the user’s objective.  Improve the system prompt,
    refine tool descriptions or add example dialogues to guide the
    model.
  \item \textbf{Low Correctness}: Indicates factual inaccuracies or
    outdated information.  Update your knowledge base, enhance
    fact‑checking logic and review external tool responses.
  \item \textbf{Low Tool Selection Accuracy}: Implies the agent chose
    the wrong tool or failed to call a tool when needed.  Adjust
    tool schemas, add explicit triggers in the system prompt, or
    refine the agent’s decision logic.
\end{itemize}

In production you might sample only a portion of sessions (e.g., 10\%)
to balance cost and coverage.  Continuous monitoring allows you to
track long‑term trends and detect regressions after updates.

\section{Summary and Next Steps}

In Lab 5 you learned how to configure online evaluations for a
production agent using AgentCore Evaluations.  You created a
configuration with built‑in evaluators, invoked the agent across
multiple scenarios and explored the resulting metrics in the
observability dashboard.  These metrics provide actionable insights
into the agent’s strengths and weaknesses and form the basis for
continuous improvement.

The final lab will build a customer‑facing web interface using
Streamlit and Cognito, connecting end users directly to the deployed
agent.  With evaluation and monitoring in place, you can confidently
roll out your agent to real customers while maintaining high quality.

